\documentclass[aspectratio=169]{beamer}

% Tema e colori
\usetheme{SimplePlus}
\usecolortheme{default}

% Pacchetti utili
\usepackage[utf8]{inputenc}
\usepackage{graphicx}
\usepackage{amsmath}
\usepackage{amssymb}

\usepackage{listings}
\usepackage{xcolor}

\lstset{
    language=C++,
    basicstyle=\ttfamily\footnotesize,
    keywordstyle=\color{blue}\bfseries,
    commentstyle=\color{green!60!black}\ttfamily,
    stringstyle=\color{red},
    breaklines=true,
    frame=none,
    tabsize=4,
    showspaces=false,
    showstringspaces=false,
    columns=flexible
}


% Informazioni sulla presentazione
\title{Clang-Tidy: Modernize}

\author{Giuseppe Di Blasi}
\institute{Università di Parma}
\date{8 Gennaio 2026}

% Logo (opzionale - decommentare e specificare il percorso)
% \logo{\includegraphics[height=1cm]{logo.png}}

\begin{document}

% Slide del titolo
\begin{frame}
\titlepage
\end{frame}

% Sezione 1
\section{Cos'è Clang-Tidy}

\begin{frame}{Cos'è Clang-Tidy}

 
  \begin{block}{Clang-Tidy}
    
  \begin{itemize}
    \item Strumento di linting basato su Clang per C++
    \item Offre un framework estensibile
    \item Individua e correggere : 
      \begin{itemize}
        \item Errori di programmazione
        \item Cattive pratiche di programmazione
        \item Problematiche di stile
      \end{itemize}
  \end{itemize}
\end{block}

  In questo seminario ci concentreremo sulla parte di modernize, quindi la rilevazione di codice legacy e 
  gli opportuni suggerimenti per adeguarlo allo standard moderno.

\end{frame}

\begin{frame}{TinyXML2}
 
  \begin{block}{TinyXML2}
    \begin{itemize}
      \item Parser per file XML scritto in C++
      \item Progetto abbastanza grande
      \item Presenta codice legacy 
    \end{itemize}
  \end{block}
   
  \begin{block}{Clang-Tidy modernize}
    \begin{itemize}
      \item Comando: \textbf{clang-tidy tinyxml2.cpp -checks="modernize-*" -{-} -std=c++17 }
    \item C++ 17 come standard target
    \item 1700 warning generati
  \end{itemize}
  \end{block}

\end{frame}

\begin{frame}
  \begin{block}{Risultato esecuzione}
    \begin{itemize}
  \item 111 modernize-use-trailing-return-type
  \item 103 modernize-use-nullptr
  \item 35 modernize-avoid-c-arrays
  \item 11 modernize-use-auto
  \item 3 modernize-use-equals-default
  \item 3 modernize-loop-convert
  \item 1 modernize-raw-string-literal
  \end{itemize}
  \end{block}
\end{frame}

% Esempi di trasformazione con clang-tidy modenrize

\section{Esempi di trasformazione}

% modernize-use-trailing-return-type

\begin{frame}[fragile]{modernize-use-trailing-return-type}
  \begin{block}{Codice legacy}
  \begin{lstlisting}
  char* StrPair::ParseText( char* p, const char* endTag, int strFlags, int* curLineNumPtr ){
    //... 
  }

   \end{lstlisting}
 \end{block}

 \begin{block}{Codice moderno}
    \begin{lstlisting}
    auto StrPair::ParseText( char* p, const char* endTag, int strFlags, int* curLineNumPtr ) -> char*{
      //... 
    }

    \end{lstlisting}
  \end{block}
\end{frame}

%%%%%%%%%%%%%%%%%%%%%%%%%%%%%%%%%%%%%%%%%%%%%%%%%%%%%%%%%%%%%%%%5

% modernize-use-nullptr

\begin{frame}[fragile]{modernize-use-nullptr}
  \begin{block}{Codice legacy}
  \begin{lstlisting}[literate={ù}{{\`u}}1]
      
      //const char* XMLUtil::GetCharacterRef(const char* p, char* value, int* length) 
      const char* q = 0; // riga 506
      
      // length_ptr sarebbe più appropriato
      if (length == 0) // riga 569
   \end{lstlisting}
 \end{block}

 \begin{block}{Codice moderno}
    \begin{lstlisting}[literate={ù}{{\`u}}1] 
    const char* q = nullptr;
    
    // length_ptr sarebbe più appropriato
    if (length == nullptr);

    \end{lstlisting}
  \end{block}

\end{frame}

%%%%%%%%%%%%%%%%%%%%%%%%%%%%%%%%%%%%%%%%%%%%%%%%%%%%%%%%%%%%%%%%%%%%%%%%%%%%%%%%%%%%%%%%%%%%%%%%%%%%%%%%%%%%%%%


% modernize-use-auto

\begin{frame}[fragile]{modernize-use-auto}
  \begin{block}{Codice legacy}
  \begin{lstlisting}
  const unsigned char* pu = reinterpret_cast<const unsigned char*>(p);

  XMLText* text = CreateUnlinkedNode<XMLText>( _textPool );
   \end{lstlisting}
 \end{block}

 \begin{block}{Codice moderno}
    \begin{lstlisting}
    auto pu = reinterpret_cast<const unsigned char*>(p);

    auto text = CreateUnlinkedNode<XMLText>( _textPool );
    \end{lstlisting}
  \end{block}

\end{frame}


%%%%%%%%%%%%%%%%%%%%%%%%%%%%%%%%%%%%%%%%%%%%%%%%%%%%%%%%%%%%%%%%%%%%%%%%%%%%%%%%%%%%%%%%%%%%%%%%%%%%%%%%%%%%%%%

% modernize-raw-string-literal
\begin{frame}[fragile]{modernize-raw-string-literal}
  \vspace{-1em}
  \begin{block}{Codice legacy}
  \begin{lstlisting}
  XMLDeclaration* XMLDocument::NewDeclaration( const char* str ){
    XMLDeclaration* dec = CreateUnlinkedNode<XMLDeclaration>( _commentPool );
    dec->SetValue( str ? str : "xml version=\"1.0\" encoding=\"UTF-8\"" );
    return dec;
  }
  \end{lstlisting}
 \end{block}

 \begin{block}{Codice moderno}
    \begin{lstlisting}
  XMLDeclaration* XMLDocument::NewDeclaration( const char* str ){
    XMLDeclaration* dec = CreateUnlinkedNode<XMLDeclaration>( _commentPool );
    dec->SetValue( str ? str : R"( xml version="1.0" encoding="UTF-8" )" );
    return dec;
  }
    \end{lstlisting}
  \end{block}

\end{frame}

%%%%%%%%%%%%%%%%%%%%%%%%%%%%%%%%%%%%%%%%%%%%%%%%%%%%%%%%%%%%%%%%%%%%%%%%%%%%%%%%%%%%%%%%%%%%%%%%%%%%%%%%%%%%%%

% modernize-avoid-c-arrays

\begin{frame}[fragile]{modernize-avoid-c-arrays}
  \vspace{-1em}
  \begin{block}{Codice legacy}
  \begin{lstlisting}
  
  static const char* TRUE_VALS[] = { "true", "True", "TRUE", 0 };
  static const char* FALSE_VALS[] = { "false", "False", "FALSE", 0 };
  \end{lstlisting}
 \end{block}

 \begin{block}{Codice moderno}
    \begin{lstlisting}
    
  static constexpr std::array<const char*, 3> TRUE_VALS = {
    "true", "True", "TRUE"
  };

  static constexpr std::array<const char*, 3> FALSE_VALS = {
    "false", "False", "FALSE"
  };
    \end{lstlisting}
  \end{block}

\end{frame}


%%%%%%%%%%%%%%%%%%%%%%%%%%%%%%%%%%%%%%%%%%%%%%%%%%%%%%%%%%%%%%%%%%%%%%%%%%%%%%%%%%%%%%%%%%%%%%%%%%%%%%%%%%%%%%%

% modernize-loop-convert

\begin{frame}[fragile]{modernize-loop-convert}
  \vspace{-1em}
  \begin{block}{Codice legacy}
  \begin{lstlisting}
  for( int i=0; i<NUM_ENTITIES; ++i ) {
    if ( entities[i].value == *q ) {
      //.. 
      Write( entities[i].pattern, entities[i].length );      
    } 
  }
  \end{lstlisting}
 \end{block}

 \begin{block}{Codice moderno}
    \begin{lstlisting}
    //si potrebbe usare const auto& per efficienza...
    for(auto entitie : entities) {
      if ( entitie.value == *q ) {
        //.. 
        Write( entitie.pattern, entitie.length );               
      } 
    }
    \end{lstlisting}
  \end{block}

\end{frame}

%%%%%%%%%%%%%%%%%%%%%%%%%%%%%%%%%%%%%%%%%%%%%%%%%%%%%%%%%%%%%%%%%%%%%%%%%%%%%%%%%%%%%%%%%%%%%%%%%%%%%%%%%%%%%%

% modernize-use-equals-default
%
\begin{frame}[fragile]{modernize-use-equals-default}
  \vspace{-1em}
  \begin{block}{Codice legacy}
  \begin{lstlisting}
     XMLComment::~XMLComment() { }
     XMLUnknown::~XMLUnknown() { }
  \end{lstlisting}
 \end{block}

 \begin{block}{Codice moderno}
    \begin{lstlisting}
    XMLComment::~XMLComment() = default;
    XMLUnknown::~XMLUnknown() = default;
    \end{lstlisting}
  \end{block}

\end{frame}


%%%%%%%%%%%%%%%%%%%%%%%%%%%%%%%%%%%%%%%%%%%%%%%%%%%%%%%%%%%%%%%%%%%%%%%%%%%%%%%%%%%%%%%%%%%%%%%%%%%%%%%%%%%%%%%

% CONCLUSIONI

\begin{frame}{Conclusioni}
  \begin{itemize}
    \item Aggiornato TinyXML2 allo standard C++ 17
    \item Facilitato l'aggiornamento allo standard moderno
    \item Migliorato la sicurezza del programma (nullptr, std::array)
    \item Migliorato la comprensione del codice con costrutti di sintassi moderna
    \item Alcuni check non fanno autofix: {modernize-avoid-c-arrays}
     \end{itemize}
\end{frame}

%%%%%%%%%%%%%%%%%%%%%%%%%%%%%5

% RIFERIMENTI 

\begin{frame}{Riferimenti}
  \begin {itemize}
\item GitHub TinyXML2: \url{https://github.com/leethomason/tinyxml2}
\item Documentazione Clang-Tidy: \url{https://clang.llvm.org/extra/clang-tidy}
\item C++ Core guidelines: \url{https://isocpp.github.io/CppCoreGuidelines/CppCoreGuidelines.html}
  \end {itemize}
\end{frame}

%%%%%%%%%%%%%%%%%%%%%%%%%%%%%%%%%%%%%%%%%%%%%%%%%%%%%%%%%%%%%%%%%%%%%%%%%%%%%%%%%%%%%%%%%%%%%%%%%%%%%%%%%%%%%%

\end{document}
